\section{Custom Protocol Specification}
\subsection{TCP control messages}

TCP control messages are UTF-8 JSON objects terminated by a newline. The receiver accumulates bytes until a complete line arrives, since TCP reads do not preserve application message boundaries. The principal message types are summarized in \cref{tab:control}.

\begin{table}[H]
\centering\small
\begin{tabularx}{\linewidth}{@{}L{3.1cm}Y@{}}
\toprule
\textbf{Message} & \textbf{Purpose} \\
\midrule
\texttt{HELLO} / \texttt{HELLO\_OK} & Agree on protocol version (currently 1). \\
\texttt{REGISTER} / \texttt{REGISTER\_OK} & Create a password account over verified TLS in account mode. \\
\texttt{LOGIN} / \texttt{LOGIN\_OK} & Authenticate and return token, roster, UDP port and session stream ID. \\
\texttt{GET\_USERS}, \texttt{USERS}, \texttt{PING} & Refresh online presence and keep a control session active. \\
\texttt{LOGOUT}, \texttt{AUTH\_REQUIRED}, \texttt{ERROR} & End a session or report authentication/protocol failures. \\
\bottomrule
\end{tabularx}
\caption{Main TCP/TLS control messages.}
\label{tab:control}
\end{table}

\subsection{UDP packet layout and encryption}

Each custom UDP frame begins with a \textbf{42-byte header} (\cref{tab:udp}). \texttt{DATA} frames append an AES-GCM encrypted payload; \texttt{ACK} frames contain the header alone. The authenticated header is used as AES-GCM associated data, so altering a DATA header invalidates decryption.

\begin{table}[H]
\centering\small
\begin{tabularx}{\linewidth}{@{}L{3.2cm}L{1.4cm}Y@{}}
\toprule
\textbf{Field} & \textbf{Bytes} & \textbf{Meaning} \\
\midrule
Magic & 4 & ASCII \texttt{COLB}; identifies this protocol. \\
Version & 1 & Protocol version, currently \texttt{1}. \\
Kind & 1 & \texttt{1 = DATA}, \texttt{2 = ACK}. \\
Sequence & 4 & Unsigned, big-endian index per peer/stream. \\
Packet ID & 16 & UUID used to match an exact retransmission/ACK. \\
Stream ID & 16 & Session-derived chat/ACK stream or file UUID. \\
Payload & variable & DATA: 12-byte nonce, ciphertext and 16-byte GCM tag; ACK: empty. \\
\bottomrule
\end{tabularx}
\caption{On-wire UDP frame structure (header = 42 bytes).}
\label{tab:udp}
\end{table}

The decrypted application payload has the format
\begin{center}
\small\texttt{[2-byte JSON length] [UTF-8 JSON metadata] [optional raw file bytes]}.
\end{center}
File bytes are not base64 encoded. The maximum UDP datagram is \textbf{1200 bytes}; file chunks contain at most \textbf{800 raw bytes}, leaving space for metadata and AES-GCM overhead. All UDP DATA payloads are encrypted; the header-only transport ACK is not encrypted.

\subsection{Application message types and reliability}

\begin{table}[H]
\centering\small
\begin{tabularx}{\linewidth}{@{}L{3.45cm}Y@{}}
\toprule
\textbf{Message} & \textbf{Use} \\
\midrule
\texttt{USERS} & Updated online-user list with monotonic revision. \\
\texttt{CHAT} & Broadcast or private message; stable \texttt{message\_id} and fixed recipient targets. \\
\texttt{CHAT\_ACK}, \texttt{CHAT\_RECEIPT} & Recipient application acceptance and sender delivery status. \\
\texttt{FILE\_START}, \texttt{FILE\_CHUNK}, \texttt{FILE\_END} & File metadata, indexed binary chunks and transfer completion marker. \\
\texttt{FILE\_RECEIPT}, \texttt{FILE\_ABORT} & Receiver hash result or transfer cancellation. \\
\texttt{ERROR} & Report invalid or failed application operations. \\
\bottomrule
\end{tabularx}
\caption{Principal UDP application messages.}
\label{tab:messages}
\end{table}

The sender retains at most \textbf{four} unacknowledged packets per stream (fixed sliding window). It retries unacknowledged packets after \textbf{0.35 s}, up to \textbf{10 retries}. Each retry retains the original sequence, packet UUID and ciphertext. The receiver buffers out-of-order DATA and suppresses duplicate delivery. Separate stream IDs prevent a file transfer from blocking the chat stream. A fixed window is \emph{not} an adaptive congestion-control algorithm; BBR/CUBIC is not implemented.
